% Standalone problem document, generated from the adjacent README.md.
% Compile with XeLaTeX. Mathematical content is maintained in Markdown.
\documentclass[11pt,a4paper]{article}
\usepackage[fontsize=10.5pt]{scrextend}
\usepackage[margin=22mm,headheight=15pt,headsep=9mm,footskip=11mm]{geometry}
\usepackage{fontspec}
\setmainfont{texgyrepagella}[Extension=.otf,UprightFont=*-regular,BoldFont=*-bold,ItalicFont=*-italic,BoldItalicFont=*-bolditalic]
\setsansfont{texgyreheros}[Extension=.otf,UprightFont=*-regular,BoldFont=*-bold,ItalicFont=*-italic,BoldItalicFont=*-bolditalic]
\setmonofont{texgyrecursor}[Extension=.otf,UprightFont=*-regular,BoldFont=*-bold,ItalicFont=*-italic,BoldItalicFont=*-bolditalic,Scale=MatchLowercase]
\usepackage{amsmath,amssymb}
\usepackage{unicode-math}
\setmathfont{texgyrepagella-math.otf}
\usepackage{xcolor,microtype,enumitem,needspace,fancyhdr}
\usepackage{bookmark}
\definecolor{ink}{HTML}{17344B}
\definecolor{accent}{HTML}{197D80}
\definecolor{muted}{HTML}{596773}
\hypersetup{colorlinks=true,linkcolor=accent,urlcolor=accent,pdftitle={MF-01 - Optimal
sign approximation with a multiplication
budget},pdfauthor={Open Problems in Numerical Linear Algebra}}
\urlstyle{same}
\setlength{\parindent}{0pt}
\setlength{\parskip}{5pt plus 1pt minus 1pt}
\linespread{1.04}
\setlength{\emergencystretch}{3em}
\widowpenalty=10000
\clubpenalty=10000
\displaywidowpenalty=10000
\allowdisplaybreaks[1]
\setlist{nosep,leftmargin=1.5em}
\providecommand{\tightlist}{\setlength{\itemsep}{0pt}\setlength{\parskip}{0pt}}
\providecommand{\pandocbounded}[1]{#1}
\pagestyle{fancy}
\fancyhf{}
\fancyhead[L]{\sffamily\footnotesize\color{muted}OPEN PROBLEMS IN NUMERICAL LINEAR ALGEBRA}
\fancyhead[R]{\sffamily\bfseries\footnotesize\color{accent}MF-01}
\fancyfoot[L]{\parbox[b]{0.9\textwidth}{\sffamily\fontsize{8}{10}\selectfont\color{muted}Editorial ratings; a bounded search cannot certify that a problem remains unsolved.\\Literature check: 2026-09-10}}
\fancyfoot[R]{\sffamily\footnotesize\color{muted}\thepage}
\renewcommand{\headrulewidth}{0.3pt}
\renewcommand{\headrule}{\hbox to\headwidth{\color{accent}\leaders\hrule height \headrulewidth\hfill}}
\makeatletter
\renewcommand{\section}{\Needspace{5\baselineskip}\@startsection{section}{1}{0pt}{12pt plus 2pt minus 2pt}{4pt}{\sffamily\bfseries\large\color{ink}}}
\renewcommand{\subsection}{\Needspace{4\baselineskip}\@startsection{subsection}{2}{0pt}{9pt plus 1pt minus 1pt}{3pt}{\sffamily\bfseries\normalsize\color{ink}}}
\makeatother
\setcounter{secnumdepth}{0}
\begin{document}
{\sffamily\small\color{muted}Matrix functions and stability\par}
\vspace{3pt}
{\raggedright\hyphenpenalty=10000\exhyphenpenalty=10000\sffamily\bfseries\fontsize{21}{24}\selectfont\color{ink}Optimal
sign approximation with a multiplication budget\par}
\vspace{7pt}
{\sffamily\small\textbf{Difficulty:} challenging\quad\textbf{Importance:} interesting
to the community\par}
{\sffamily\small\bfseries\color{accent}Status: Open\par}
\vspace{4pt}
{\color{accent}\hrule height 0.8pt}
\vspace{6pt}
\textbf{Rating rationale:} Challenging because multiplication complexity
and uniform approximation must be optimized together; community impact
comes from sign, projector, and polar computations.

\subsection{Problem statement}\label{problem-statement}

For \(m\in\mathbb N_0\) and \(0<\delta<1\), put

\[
I_\delta=[-1,-\delta]\cup[\delta,1].
\]

Let \(\mathcal P_m\) consist of real polynomials computed from \(1,x\)
by straight-line programs using at most \(m\) nonscalar multiplications;
real linear combinations cost nothing. Define

\[
E_m(\delta)=\inf_{p\in\mathcal P_m}
\max_{x\in I_\delta}|p(x)-\operatorname{sign}(x)|.
\]

Determine matching asymptotic upper and lower bounds for
\(E_m(\delta)\), with the dependence on both \(m\) and \(\delta\)
explicit. Infima avoid assuming attainment. The arithmetic model
concerns a single polynomial identity valid for matrices of every size.

\subsection{References and status
evidence}\label{references-and-status-evidence}

Amsel et al., \href{https://arxiv.org/html/2602.05394v3}{Linear Systems
and Eigenvalue Problems: Open Questions from a Simons Workshop}, §6.3,
Problem 6.3 (2026). Rubensson, Jarlebring and Lorentzon,
\href{https://arxiv.org/html/2606.24701v1}{Recursive expansion of the
matrix step function using polynomials of degree eight}, §7 (June 2026),
still identifies optimal recursive expansion as open. Its particular
algorithm does not establish the unrestricted optimum above.

\subsection{Audit --- 2026-09-10}\label{audit-2026-09-10}

Rechecked the
\href{https://arxiv.org/html/2602.05394v3\#S6.SS3}{workshop's Problem
6.3} and the \href{https://arxiv.org/html/2606.24701v1}{degree-eight
follow-up, §7}. Neither determines the unrestricted optimum. Searches
for optimal sign approximation and recursive-expansion follow-ups found
no resolution; this is a bounded literature check.
\end{document}
